\section{The \texttt{Slit\_N} McXtrace Component}
Release: McXtrace 0.1

Rectangular/circular slit, duplicated

\subsection*{Identification}
\begin{itemize}
  \item \textbf{Author:} Erik Knudsen
  \item \textbf{Origin:} Risoe
  \item \textbf{Date:} June 16, 2009
\end{itemize}

\subsection*{Description}
Based on Slit-comp by Kim Lefmann and Henrik Roennow A simple rectangular or circular slit. You may either specify the radius (circular shape), which takes precedence, or the rectangular bounds. The slits are separated with 'd', and arranged along X. No transmission around the slit is allowed. If cutting option is used, low-weight x-rays are ABSORBED

Example: Slit\_N(xwidth=0.01, yheight=0.01, d=0.02) Slit\_N(radius=0.01, cut=1e-10, d=0.03)

\subsection*{Input parameters}
Parameters in \textbf{boldface} are required; the others are optional.

\begin{longtable}{p{0.22\textwidth}p{0.12\textwidth}p{0.46\textwidth}p{0.14\textwidth}}
\toprule
\textbf{Name} & \textbf{Unit} & \textbf{Description} & \textbf{Default} \\
\midrule
\endhead
radius & m & Radius of slit in the z=0 plane, centered at Origo & 0 \\
cut &  & Lower limit for allowed weight (1) & 0 \\
xwidth & m & Width of slit. Overrides xmin,xmax. & 0 \\
yheight & m & Height of slit. Overrides ymin,ymax. & 0 \\
N & 1 & Number of slit openings along X & 2 \\
d & m & Separation of slits & 0 \\
\bottomrule
\end{longtable}

\subsection*{Links}
\begin{itemize}
  \item Component source code found in file \texttt{Slit\_N.comp}.
\end{itemize}
\IfFileExists{optics/Slit_N_static.tex}{\input{optics/Slit_N_static.tex}}{}